% !TEX root = InformatiqueMP2I.tex

\documentclass[a4paper,12pt,titlepage]{article}
\usepackage[utf8]{inputenc}
\usepackage[french]{babel}
\usepackage[T1]{fontenc}
\usepackage[margin=2.5cm]{geometry}
\usepackage{amsmath, amsfonts, amssymb}
\usepackage{xcolor}
\usepackage[most]{tcolorbox}
\usepackage[finalizecache=true]{minted}
\usepackage{tabularx}
\usepackage{hyperref}
\usepackage{graphics}
\usemintedstyle{perso}
\renewcommand{\tabularxcolumn}[1]{>{\arraybackslash}m{#1}}

\renewcommand{\c}{\subsection[C]{C\hfill\protect\includegraphics[height=11pt]{C}}}
\newcommand{\ocaml}{\subsection[OCaml]{OCaml\hfill\protect\includegraphics[height=11pt]{Ocaml}}}

\let\oldsection\section
\renewcommand\section{\clearpage\oldsection}

\newcolumntype{A}{@{}>{\raggedright\arraybackslash\hspace{0pt}}p{.2\hsize}@{}}
\newcolumntype{B}{@{}>{\raggedright\arraybackslash\hspace{0pt}}p{.8\hsize}X@{}}

\let\oldleft\left
\renewcommand{\left}{\mathopen{}\mathclose\bgroup\oldleft}
\let\oldright\right
\renewcommand{\right}{\aftergroup\egroup\oldright}
\renewcommand{\O}[1]{\ensuremath{O\left(#1\right)}}
\newcommand{\Th}[1]{\ensuremath{\varTheta\left(#1\right)}}

\definecolor{kywd}{RGB}{17,96,132}
\definecolor{fnc}{RGB}{160,160,0}
\definecolor{var}{RGB}{0,40,186}
\newcommand{\kywd}[1]{\texttt{\color{kywd}#1}}
\newcommand{\vr}[1]{\mathtt{\color{var}#1}}

\setlength{\parindent}{0pt}
\newtcbtheorem{fun}{Fonction disponible }{breakable,colback=green!5,colframe=green!75!black,fonttitle=\bfseries}{func}
\newtcbtheorem{imp}{Implémentation }{breakable,colback=orange!5,colframe=orange!50!red,fonttitle=\bfseries}{impl}
\newtcbtheorem{tp}{Définition d'un type }{breakable,colback=violet!5,colframe=violet!75!black,fonttitle=\bfseries}{type}
\newtcbtheorem{fnc}{Fonction à implémenter }{breakable,colback=blue!5,colframe=blue!75!white,fonttitle=\bfseries}{func}
\newtcbtheorem{df}{Définition }{breakable,colback=teal!5,colframe=teal!75!white,fonttitle=\bfseries}{defi}

\newcommand{\info}[4][]{\def\argI{}\ifx&#2&\else\def\argI{Signature&\mintinline{ocaml}|#2|\\}\fi\def\argII{}\ifx&#3&\else\def\argII{Description&#3\\}\fi\begin{tabularx}{\textwidth}{AB}\argI\argII Complexité&#4#1\end{tabularx}}
\newcommand{\fn}[5]{\begin{fun*}{\mintinline{#1}|#2|}\info{#3}{#4}{#5}\end{fun*}}

\hypersetup{hidelinks,linktoc=all}
\setcounter{secnumdepth}{5}
\setcounter{tocdepth}{5}
\hypersetup{pdftitle=Informatique MP2I}
\title{Informatique MP2I}
\author{}
\date{}

\begin{document}
\pagenumbering{gobble}
\maketitle
\null\newpage
\tableofcontents
\null\newpage
\null\newpage
\pagenumbering{arabic}
\section{Listes chaînées}
%%%%%%%%%%%%%%%%%%%%%%%%%%%%%%%%%%%%%%%%%%%%%%%%%%%%%%%%%%%%%%%%%%%%%%%%%%%%%%%%%%%%%%%%%%%%%%%%%%%
\c
\begin{tp*}{}
\begin{minted}{c}
struct cell {
    int value;
    struct cell* next;
};
typedef struct cell int_list;
/*
    Cette définition de listes chaînées ainsi que la majorité des
    fonctions à suivre s'adaptent également pour les autres types
*/
\end{minted}
\end{tp*}
\begin{fnc*}{premier élément d'une liste}
\begin{minted}{c}
int fst(int_list* lst) {
    return int_list->value;
}
\end{minted}
\info{}{Retourne le premier élément d'une liste pointée par \mintinline{c}|lst|}{\Th{1}}
\end{fnc*}
\begin{fnc*}{dernier élément d'une liste}
\begin{minted}{c}
int last(int_list* lst) {
    while (lst->nest != NULL) {lst = lst->next;}
    return int_list->value;
}
\end{minted}
\info{}{Retourne le dernier élément d'une liste pointée par \mintinline{c}|lst|}{\Th{n}}
\end{fnc*}
\begin{fnc*}{longueur d'une liste}
\begin{minted}{c}
int length(int_list* lst) {
    if (lst == NULL) {
        return 0;
    }
    return 1 + length(lst->next);
}
\end{minted}
\info{}{Renvoie la longueur de la liste}{\Th{n}}
\end{fnc*}
\begin{fnc*}{ajout d'un élément à gauche d'une liste}
\begin{minted}{c}
void add_l(int elem, int_list* lst) {
    int_list* new_p = (int_list*)malloc(sizeof(int_list));
    new_p->next = lst;
    new_p->value = elem;
}
\end{minted}
\info{}{Ajoute un élément \mintinline{c}|elem| à gauche de la liste pointée par \mintinline{c}|lst|}{\Th{1}}
\end{fnc*}
\begin{fnc*}{ajout d'un élément à droite d'une liste}
\begin{minted}{c}
void add_r(int elem, int_list* lst) {
    while (lst->next != NULL) {lst = lst->next;}
    int_list* new_p = (int_list*)malloc(sizeof(int_list));
    new_p->next = NULL;
    new_p->value = elem;
    lst->next = new_p;
}
\end{minted}
\info{}{Ajoute un élément \mintinline{c}|elem| à droite de la liste pointée par \mintinline{c}|lst|}{\Th{n}}
\end{fnc*}
\begin{fnc*}{appartenance à une liste}
\begin{minted}{c}
bool mem(int elem, int_list* lst) {
    if (lst->value == elem) {return true;}
    if (lst->next != NULL) {
        return mem(elem, lst->next);
    }
}
\end{minted}
\info{}{Vérifie si un élément \mintinline{c}|elem| est dans la liste pointée par \mintinline{c}|lst|}{\O{n}}
\end{fnc*}
\begin{fnc*}{retournement d'une liste}
\begin{minted}{c}
voir rev(int_list* lst) {
    int_list* rev_p = NULL;
    for (int_list* ptr = p_liste; ptr != NULL; ptr = ptr->next) {
        liste_int* tmp = (int_list*)malloc(sizeof(int_list));
        tmp->value = ptr->valeur;
        tmp->next = rev_p;
        rev_p = tmp;
    }
    return rev_p;
}
\end{minted}
\info{}{Retourne les éléments de la liste pointée par \mintinline{c}|lst|}{\Th{n}}
\end{fnc*}
\begin{fnc*}{suppression d'une liste}
\begin{minted}{c}
void del(int_list* lst){
    while (lst->next != NULL) {
        list_int* tmp = lst;
        lst = tmp->next;
        free(tmp);
    }
}
\end{minted}
\info{}{Supprime une liste pointée par \mintinline{c}|lst|}{\Th{n}}
\end{fnc*}
%%%%%%%%%%%%%%%%%%%%%%%%%%%%%%%%%%%%%%%%%%%%%%%%%%%%%%%%%%%%%%%%%%%%%%%%%%%%%%%%%%%%%%%%%%%%%%%%%%%
\ocaml
\subsubsection{Type complexe}
\subsubsection{Le module \kywd{List}}
\fn{ocaml}{List.length}{'a list -> int}{Renvoie la longueur de la liste}{\Th{n}}
\begin{imp*}{}
\begin{minted}{ocaml}
let rec length = function
    | [] -> 0
    | h::t -> 1 + length t;;
\end{minted}
\info{'a list -> int}{Renvoie la longueur de la liste}{\Th{n}}
\end{imp*}
\fn{ocaml}{List.iter}{('a -> unit) -> 'a list -> unit}{Applique une fonction à tous les éléments de la liste}{\Th{n}}
\begin{imp*}{}
\begin{minted}{ocaml}
let rec iter f = function
    | [] -> ()
    | h::t -> f h; iter f t;;
\end{minted}
\info{('a -> unit) -> 'a list -> unit}{Applique une fonction à tous les éléments de la liste}{\Th{n}}
\end{imp*}
\fn{ocaml}{List.fold_left}{('a -> 'b -> 'a) -> 'a -> 'b list -> 'a}{Applique une fonction successivement à un élément de la liste et au résultat de l'itération précédente en partant de la fin de la liste}{\Th{n}}
\begin{imp*}{}
\begin{minted}{ocaml}
let rec fold_left f acc = function
    | [] -> acc
    | h::t -> fold_left f (f acc h) t;;
\end{minted}
\info{('a -> 'b -> 'a) -> 'a -> 'b list -> 'a}{Applique une fonction successivement à un élément de la liste et au résultat de l'itération précédente en partant de la fin de la liste}{\Th{n}}
\end{imp*}
\fn{ocaml}{List.map}{('a -> 'b) -> 'a list -> 'b list}{Applique une fonction à tous les éléments de la liste et renvoie une nouvelle liste avec les résultats}{\Th{n}}
\begin{imp*}{}
\begin{minted}{ocaml}
let rec map f = function
    | [] -> []
    | h::t -> (f h)::(map f t);;
\end{minted}
\info{('a -> 'b) -> 'a list -> 'b list}{Applique une fonction à tous les éléments de la liste et renvoie une nouvelle liste avec les résultats}{\Th{n}}
\end{imp*}
\fn{ocaml}{List.rev}{'a list -> 'a list}{Inverse l'ordre des éléments d'une liste}{\Th{n}}
\begin{imp*}{}
\begin{minted}{ocaml}
let rec rev l =
    let rec aux acc = function
        | [] -> acc
        | h::t -> aux (h::acc) t
    in aux [] l;;
\end{minted}
\info{'a list -> 'a list}{Inverse l'ordre des éléments d'une liste}{\Th{n}}
\end{imp*}
\fn{ocaml}{List.mem}{'a -> 'a list -> bool}{Vérifie si un élément appartient à une liste}{\O{n}}
\begin{imp*}{}
\begin{minted}{ocaml}
let rec mem elem = function
    | [] -> false
    | h::t -> elem = h || mem elem t;;
\end{minted}
\info{'a -> 'a list -> bool}{Vérifie si un élément appartient à une liste}{\O{n}}
\end{imp*}
\section{Tris}
\begin{df*}{Tri en place}
Tri qui déplace les éléments à l'intérieur du tableau au lieu d'en créer un nouveau.
\end{df*}
\begin{df*}{Tri stable}
Tri qui préserve l'ordre des éléments (si \mintinline{c}|a| apparaît avant \mintinline{c}|b| dans le tableau initial, alors \mintinline{c}|a| apparaît avant \mintinline{c}|b| dans le tableau trié).
\end{df*}
%%%%%%%%%%%%%%%%%%%%%%%%%%%%%%%%%%%%%%%%%%%%%%%%%%%%%%%%%%%%%%%%%%%%%%%%%%%%%%%%%%%%%%%%%%%%%%%%%%%
\subsection{Différents algorithmes de tris}
\begin{df*}{Tri par sélection}
Tri qui sélectionne le plus petit élement du tableau non trié et l'insère au début de ce dernier.
\end{df*}
\begin{df*}{Tri par insertion}
Tri qui parcourt le tableau de gauche à droite, insérant chaque élément à sa place dans la partie triée du tableau.
\end{df*}
\begin{df*}{Tri à bulles}
Tri qui parcourt le tableau de gauche à droite en comparant chaque élément avec son voisin, et échange les deux éléments si le voisin est plus petit.
\end{df*}
\begin{pt*}{Tri à bulles}
Le tri à bulles est stable.
\end{pt*}
\begin{df*}{Tri rapide}
Tri qui sélectionne un pivot et partitionne le tableau en deux sous-tableaux: les éléments plus petits que le pivot et ceux plus grands. Cette opération est répétée récursivement sur chaque sous-tableau.
\end{df*}
\begin{df*}{Tri mixte}
Tri qui utilise plusieurs algorithmes de tri pour avoir une meilleure complexité.
\end{df*}
\c
\begin{imp*}{échange de deux élements dans un tableau}
\begin{minted}{c}
void array_swap(int* tab, int i, int j){
    int tmp = tab[i];
    tab[i] = tab[j];
    tab[j] = tmp;
}
\end{minted}
\info{}{Échange les éléments \mintinline{c}|i| et \mintinline{c}|j| du tableau \mintinline{c}|tab|}{\Th{1}}
\end{imp*}
\begin{imp*}{tri sélection}
\begin{minted}{c}
void tri_selection(int* tab, int n) {
    for (int i = 0; i < n-1; ++i) {
        int min = tab[i];
        int idx_min = i;
        for (int j = i+1; j < n; ++j) {
            if (tab[j] < min) {
                min = tab[j];
                idx_min = j;
            }
        }
        array_swap(tab, i, idx_min);
    }
}
\end{minted}
\info{}{Tri le tableau \mintinline{c}|tab| avec un tri sélection}{\Th{n^2}}
\end{imp*}
\begin{imp*}{tri insertion}
\begin{minted}{c}
void tri_insertion(int* tab, int n) {
    for (int i = 1; i < n; ++i) {
        for (int j = i; j > 0 && tab[j] < tab[j-1]; --j) {
            array_swap(tab, j, j-1);
        }
    }
}
\end{minted}
\info{}{Tri le tableau \mintinline{c}|tab| avec un tri insertion}{\Th{n^2}}
\end{imp*}
\begin{imp*}{tri à bulles}
\begin{minted}{c}
void tri_bulle(int* tab, int n) {
    for (int i = 0; i < n; ++i) {
        for (int j = n-1; j > i; --j) {
            if (tab[j] < tab[j-1]) {
                array_swap(tab, j, j-1);
            }
        }
    }
}
\end{minted}
\info{}{Tri le tableau \mintinline{c}|tab| avec un tri à bulles}{\Th{n^2}}
\end{imp*}
\begin{imp*}{partitionnement d'un tableau pour le tri rapide}
\begin{minted}{c}
int partition(int* tab, int n) {
    int pivot = tab[0];
    int d = n-1;
    int g = 1;
    while (d > g-1) {
        if (tab[g] < pivot) {++g;}
        else {array_swap(tab, g, d); --d;}
    }
    array_swap(tab, 0, g-1);
    return g;
}
\end{minted}
\info{int* tab, int n}{Partitionne le tableau \mintinline{c}|tab| de taille \mintinline{c}|n| pour le tri rapide}{\Th{n}}
\end{imp*}
\begin{imp*}{tri rapide}
\begin{minted}{c}
void tri_rapide(int* tab, int n) {
    int p = partition(tab, n);
    if (p > 1) {tri_rapide(tab, p);}
    if (n-p > 1) {tri_rapide(&tab[p], n-p);}
}
\end{minted}
\info{}{Trie le tableau \mintinline{c}|tab| avec un tri rapide}{\Th{n \log n} (en moyenne), \Th{n^2} (dans le pire des cas)}
\end{imp*}
\begin{imp*}{tri rapide}
\begin{minted}{c}
void tri_mixte(int* tab, int n) {
    if (n < 10) {
        tri_bulle(tab,n);
    } else {
        int p = partition(tab, n);
        if (p > 1) {tri_mixte(tab, p);}
        if (n-p > 1) {tri_mixte(&tab[p], n-p);}
    }
}
\end{minted}
\info{}{Trie le tableau \mintinline{c}|tab| avec un tri rapide ou un tri à bulles si la liste est petite}{\Th{n \log n} (en moyenne), \Th{n^2} (dans le pire des cas)}
\end{imp*}
\ocaml
\begin{imp*}{minimum d'une liste}
\begin{minted}{ocaml}
let min = function
    | [] -> failwith "Liste vide"
    | t::q ->
        let rec lst_min t = function
            | [] -> t;
            | t2::q2 when t < t2 -> lst_min t q2
            | t2::q2 -> lst_min t2 q2
        in lst_min t q;;
\end{minted}
\info{'a list -> 'a}{Renvoie le minimum de la liste fournie en paramètre}{\Th{n}}
\end{imp*}
\begin{imp*}{retire un élément d'une liste}
\begin{minted}{ocaml}
let rec retire elem = function
    | [] -> failwith "Liste vide"
    | t::q when t = elem -> q
    | t::q -> t::retire elem q;;
\end{minted}
\info{'a list -> 'a list}{Retire l'élément \mintinline{ocaml}|elem| de la liste fournie en paramètre}{\Th{n}}
\end{imp*}
\begin{imp*}{tri sélection}
\begin{minted}{ocaml}
let rec tri_selection = function
    | [] -> []
    | lst -> let m = min lst in m::tri_selection (retire m lst);;
\end{minted}
\info{'a list -> 'a list}{Trie le tableau avec un tri sélection}{\Th{n^2}}
\end{imp*}
\begin{imp*}{Tri insertion}
\begin{minted}{ocaml}
let rec tri_insertion = function
    | [] -> []
    | t::q ->
        let rec insere elem = function
            | [] -> [elem]
            | t::q when t < elem -> t::insere elem q
            | lst -> elem::lst
        in insere t (tri_insertion q);;
\end{minted}
\info{'a list -> 'a list}{Trie le tableau avec un tri insertion}{\Th{n^2}}
\end{imp*}
\begin{imp*}{tri bulle}
\begin{minted}{ocaml}
let rec tri_bulle = function
    | [] -> []
    | lst ->
        let rec bulle = function
            | [] -> []
            | t::[] -> [t]
            | t::q ->
                let b = bulle q in
                    if List.hd b < t
                    then List.hd b::t::(List.tl b)
                    else t::b
        in let b = bulle lst
        in List.hd b::tri_bulle (List.tl b);;
\end{minted}
\info{'a list -> 'a list}{Trie le tableau avec un tri bulle}{\Th{n^2}}
\end{imp*}
\begin{imp*}{joint deux listes}
\begin{minted}{ocaml}
let rec join lst = function
    | [] -> lst
    | t::q -> t::join lst q
\end{minted}
\info{'a list -> 'a list -> 'a list}{Joint une liste devant la liste \mintinline{ocaml}|lst|}{\Th{n} ($n$ est la longueur de la deuxième liste)}
\end{imp*}
\begin{imp*}{partitionnement d'une liste}
\begin{minted}{ocaml}
let rec partition p left right = function
    | [] -> (left,right)
    | t::q when t < p -> partition p (t::left) right q
    | t::q -> partition p left (t::right) q;;
\end{minted}
\info{'a -> 'a list -> 'a list -> 'a list|\newline\mintinline{ocaml}|    -> 'a list * 'a list}{Partitionne les éléments de la liste donnée en paramètre\newline\mintinline{ocaml}|left| et \mintinline{ocaml}|right| servent d'accumulateurs pour les éléments plus petits et plus grand que le pivot \mintinline{ocaml}|p|}{\Th{n}}
\end{imp*}
\begin{imp*}{tri rapide}
\begin{minted}{ocaml}
let rec tri_rapide = function
    | [] -> []
    | t::q ->
        let (l,r) = partition t [] [] q in
        join (t::tri_rapide r) (tri_rapide l);;
\end{minted}
\info{'a list -> 'a list}{Trie le tableau avec un tri rapide}{\Th{n \log n} (en moyenne), \Th{n^2} (dans le pire des cas)}
\end{imp*}

\section{\emph{Tablistes}}
\begin{df*}{\emph{Tabliste}}
Structure de données qui permet de stocker un sous-ensemble $\mathcal{N}$ fini de $\mathbb{N}$ (de la forme $\left[\!\left[0,n-1\right]\!\right]$), et d'effectuer des opérations sur celle-ci en \Th{1}.
\end{df*}
%%%%%%%%%%%%%%%%%%%%%%%%%%%%%%%%%%%%%%%%%%%%%%%%%%%%%%%%%%%%%%%%%%%%%%%%%%%%%%%%%%%%%%%%%%%%%%%%%%%
\c
\begin{tp*}{}
\begin{minted}{c}
typedef struct {int* pos; int* values; int size;} Tablist;
\end{minted}
La champ \mintinline{c}|values| correspond à une liste des entiers de $\mathcal{N}$ sous la forme d'un tableau de taille $n$ dont les \mintinline{c}|size| premières cases contiennent les éléments présents dans $\mathcal{N}$.\\Le champ \mintinline{c}|pos| est un tableau de taille $n$ tel que: si $\vr{k}\in\mathcal{N}$, \mintinline{c}|pos[k]| contient la position de \mintinline{c}|k| dans la liste \mintinline{c}|values| et une valeur quelconque sinon.
\end{tp*}
\begin{fnc*}{création d'une \emph{tabliste}}
\begin{minted}{c}
Tablist init(int n) {
    Tablist t = {
        .pos = malloc(n*sizeof(int)),
        .values = malloc(n*sizeof(int)),
        .size = 0
    };
    for (int i = 0; i < n; ++i) {
        t.pos[i] = 0; // nécessaire à cause des nouvelles normes
        t.values[i] = 0; // facultatif
    }
    return t;
}
\end{minted}
\info{}{Créé une \emph{tabliste} vide}{\Th{n}}
\end{fnc*}
\begin{fnc*}{appartenance à une \emph{tabliste}}
\begin{minted}{c}
bool mem(Tablist t, int k) {
    int p = t.pos[k];
    return (p > 0 && p <= t.size && t.values[p]==k);
}
\end{minted}
\info{}{Vérifie si un élement \mintinline{c}|k| appartient à la \emph{tabliste} \mintinline{c}|t|}{\Th{1}}
\end{fnc*}
\begin{fnc*}{ajout à une \emph{tabliste}}
\begin{minted}{c}
void add(Tablist* ptr_t, int k) {
    if (!mem(*ptr_t, k)) {
        ptr_t->values[ptr_t->size] = k;
        ptr_t->pos[k] = ptr_t->size;
        ++ptr_t->size;
    };
}
\end{minted}
\info{}{Ajoute un élément \mintinline{c}|k| à la \emph{tabliste} pointée par \mintinline{c}|ptr_t|}{\Th{1}}
\end{fnc*}
\begin{fnc*}{supression d'une \emph{tabliste}}
\begin{minted}{c}
void t_remove(Tablist* ptr_t, int k) {
    if (mem(*ptr_t,k)) {
        int i = ptr_t->values[ptr_t->size-1];
        int p = ptr_t->pos[k];
        ptr_t->values[p] = i;
        ptr_t->pos[i] = p;
        --ptr_t->size;
    }
}
\end{minted}
\info{}{Supprime l'élément \mintinline{c}|k| de la \emph{tabliste} pointée par \mintinline{c}|ptr_t|}{\Th{1}}
\end{fnc*}
\begin{fnc*}{affichage d'une \emph{tabliste}}
\begin{minted}{c}
void print(Tablist t) {
    for (int i = 0; i < t.size-1; ++i) {
        printf("%d, ",t.values[i]);
    }
    printf("%d\n",t.values[t.size-1]);
}
\end{minted}
\info{}{Affiche les éléments de la \emph{tabliste} \mintinline{c}|t| dans laquelle ils ont été ajoutés à celle-ci}{\Th{\left|\mathcal{N}\right|}}
\end{fnc*}
\begin{fnc*}{vidage d'une \emph{tabliste}}
\begin{minted}{c}
void empty(Tablist* ptr_t) {
    ptr_t->size = 0;
}
\end{minted}
\info{}{Supprime les éléments de la \emph{tabliste} pointée par \mintinline{c}|ptr_t| ($\mathcal{N}=\varnothing$)}{\Th{1}}
\end{fnc*}
%%%%%%%%%%%%%%%%%%%%%%%%%%%%%%%%%%%%%%%%%%%%%%%%%%%%%%%%%%%%%%%%%%%%%%%%%%%%%%%%%%%%%%%%%%%%%%%%%%%
\ocaml
\begin{tp*}{}
\begin{minted}{ocaml}
type tablist = {
    pos : int array;
    values : int array;
    mutable size : int
};;
\end{minted}
La champ \mintinline{c}|values| correspond à une liste des entiers de $\mathcal{N}$ sous la forme d'un tableau de taille $n$ dont les \mintinline{c}|size| premières cases contiennent les éléments présents dans $\mathcal{N}$.\\Le champ \mintinline{c}|pos| est un tableau de taille $n$ tel que: si $\vr{k}\in\mathcal{N}$, \mintinline{c}|pos[k]| contient la position de \mintinline{c}|k| dans la liste \mintinline{c}|values| et une valeur quelconque sinon.
\end{tp*}
\begin{fnc*}{création d'une \emph{tabliste}}
\begin{minted}{ocaml}
let init n = {
    pos = Array.make n 0;
    values = Array.make n 0;
    size = 0
};;
\end{minted}
\info{int -> tablist}{Créé une \emph{tabliste} vide}{\Th{n}}
\end{fnc*}
\begin{fnc*}{appartenance à une \emph{tabliste}}
\begin{minted}{ocaml}
let mem t k =
    let p = t.pos.(k) in
    p > 0 && p < t.size && t.values.(p) = k;;
\end{minted}
\info{tablist -> int -> bool}{Vérifie si un élement \mintinline{c}|k| appartient à la \emph{tabliste} \mintinline{c}|t|}{\Th{1}}
\end{fnc*}
\begin{fnc*}{ajout à une \emph{tabliste}}
\begin{minted}{ocaml}
let add t k =
    if not (mem t k) then
        begin
            t.values.(t.size) <- k;
            t.pos.(k) <- t.size;
            t.size <- t.size + 1;
        end;
    ;;
\end{minted}
\info{tablist -> int -> unit}{Ajoute un élément \mintinline{c}|k| à la \emph{tabliste} pointée par \mintinline{c}|ptr_t|}{\Th{1}}
\end{fnc*}
\begin{fnc*}{supression d'une \emph{tabliste}}
\begin{minted}{ocaml}
let remove t k =
    if mem t k then
        begin
            let i = t.values.(t.size - 1)
            and p = t.pos.(k) in
            t.values.(p) <- i;
            t.pos.(i) <- p;
            t.size <- t.size - 1;
        end;
    ;;
\end{minted}
\info{tablist -> int -> unit}{Supprime l'élément \mintinline{c}|k| de la \emph{tabliste} pointée par \mintinline{c}|ptr_t|}{\Th{1}}
\end{fnc*}
\begin{fnc*}{affichage d'une \emph{tabliste}}
\begin{minted}{ocaml}
let print t =
    for i = 0 to t.size - 2 do
        print_int t.values.(i);
        print_string ", ";
    done;
    print_int t.values.(t.size - 1);
    print_newline ();;
\end{minted}
\info{tablist -> unit}{Affiche les éléments de la \emph{tabliste} \mintinline{c}|t| dans laquelle ils ont été ajoutés à celle-ci}{\Th{\left|\mathcal{N}\right|}}
\end{fnc*}
\begin{fnc*}{vidage d'une \emph{tabliste}}
\begin{minted}{ocaml}
let empty t = t.size <- 0;;
\end{minted}
\info{tablist -> unit}{Supprime les éléments de la \emph{tabliste} pointée par \mintinline{c}|ptr_t| ($\mathcal{N}=\varnothing$)}{\Th{1}}
\end{fnc*}
\pagenumbering{gobble}
\null\newpage
\end{document}